\section{Exact finite algebra and the public detector}\label{sec:exact-algebra}

This section makes the finite layer of Section~4 self-contained.  The
geometric binding of the same data to the Johnson replacement collar is
Hypothesis~\ref{hyp:P1}, which is \Open; the integer $2624$ itself is
independent of that binding and is the value of the frozen Burau
computation.

\subsection{Erratum to the 2 September public draft}

The 2 September public draft reported $D_3(v_T)=-59072$.  That value mixed
two endpoint coordinate systems: the vector $u$ used the THXY endpoint
indices, whereas the covector $\ell$ used the collar indices that also label
the cabled Artin word.  The two tables order the $m_2$ wickets in opposite
directions.  In the collar table, the same physical cup has endpoints $2$
and $87$, so
\[
  u=e_2-e_{87},\qquad \ell=e_{87}^*-e_2^*.
\]
With both objects in that one table, the exact calculation gives
$[\varepsilon^3]\ell(\rho(W)-I)u=-328$ and hence the cubic coefficient is
$(-2)^3(-328)=2624$.  The nonvanishing conclusion is unchanged.
Placing both $u$ and $\ell$ in the THXY coordinates while leaving the word
in collar coordinates gives $-2496$; this is not a coordinate change of the
whole pairing, because the operator is not transported, and it is not a
second geometric class (Appendix~\ref{app:legality}).  The endpoint
transport program reproduces the withdrawn numbers as illegal partial
transports: THXY cup with collar cap, $-59072$; THXY cup and cap with the
untransported word, $-2496$; frozen cut-slot cup with collar cap,
$-55104$; cut-slot cup and cap with the untransported word, $-115456$.
Simultaneous transport of the word, the cup and the cap gives $2624$ in
every coordinate system, and the derived constant terms are exactly
\eqref{eq:u} and \eqref{eq:ell}.  The value $2624$ belongs to the
currently frozen collar data; if the Johnson collar word changes, the
numerical value must be recomputed.

\subsection{Integral inverses and the long topology route}

Besides \eqref{eq:determinants}, define
\begin{equation}\label{eq:A-minus-I}
A-I=
\begin{pmatrix}
-1&269&1240\\
0&40&189\\
1&0&31
\end{pmatrix}.
\end{equation}
Expansion along the first row gives $\det(A-I)=1$.  An integral inverse is
\begin{equation}\label{eq:B-inverse}
(A-I)^{-1}_{\Z}=
\begin{pmatrix}
1240&-8339&1241\\
189&-1271&189\\
-40&269&-40
\end{pmatrix}.
\end{equation}
Since $\det A=1$, direct inversion gives
\begin{equation}\label{eq:H2-map}
C=(A^{-1})^T-I=
\begin{pmatrix}
1311&189&-41\\
-8608&-1241&269\\
1&0&-1
\end{pmatrix},
\end{equation}
with integral inverse
\begin{equation}\label{eq:H2-inverse}
C^{-1}_{\Z}=
\begin{pmatrix}
-1241&-189&40\\
8339&1270&-269\\
-1241&-189&39
\end{pmatrix}.
\end{equation}
These matrices satisfy $\det A=\det(A-I)=1$.  The sphere coordinate matrix in
\eqref{eq:spherecols} is obtained by negating the columns of
\eqref{eq:H2-map} and also has determinant $1$.

\begin{remark}[Long topology route]\label{rem:long-topology-route}
Surjectivity of $A-I$ kills the natural coinvariant presentation of the
fundamental group after surgery; bijectivity of $A-I$ and $(A^{-1})^T-I$
kills the intermediate homology in the Wang and Mayer--Vietoris sequences.
Poincar\'e duality then gives the homology profile of $S^4$.  Hurewicz and
Whitehead promote simple connectivity plus that homology profile to a
homotopy equivalence with $S^4$.  This paper uses Iwaki
Proposition~2.1 on the identified surgery manifold instead.
\end{remark}

\begin{corollary}[Homotopy-sphere conclusion]
\label{cor:homotopy-sphere}
The standard Cappell--Shaneson manifold $\Sigma_A^0$ is a homotopy
$4$-sphere.  If the Johnson CS handle picture $X_J$ is identified with
$\Sigma_A^0$ (P3/E13, \Open), the same holds for $X_J$.
\end{corollary}

\begin{proof}
By \eqref{eq:determinants}, $A\in\mathrm{SL}(3,\Z)$ and $\det(A-I)=1$, so
Iwaki Proposition~2.1 \cite{Iwaki2024} applies to $\Sigma_A^0$.  The
identification $X_J\cong\Sigma_A^0$ depends on Hypothesis~\ref{hyp:P0}.
\end{proof}

\subsection{Coefficient traces and the selected Hattori class}

\begin{definition}[Coefficient trace]\label{def:coefficient-HH0}
For a small $\Q$-linear category $\mathcal C$ and coefficient bimodule $M$,
the coefficient zeroth Hochschild homology is
\[
\HH(\mathcal C;M)=
\left(\bigoplus_{T}M(T,T)\right)
\Big/
\Span_{\Q}\{f\cdot m-m\cdot f\},
\]
where $f:S\to T$ and $m\in M(T,S)$ range over every cyclically composable
pair.
\end{definition}

\begin{lemma}[Descent through the coefficient trace]\label{lem:HH0-descent}
Let $r_T:M(T,T)\to V$ satisfy $r_S(f\cdot m)=r_T(m\cdot f)$ for all
cyclically composable $f$ and $m$.  Then there is a unique linear map
$\bar r:\HH(\mathcal C;M)\to V$ whose restriction to $M(T,T)$ is $r_T$.
\end{lemma}

\begin{proof}
The sum of the $r_T$ kills every relation generator $f\cdot m-m\cdot f$ and
therefore factors uniquely through the quotient.
\end{proof}

For the Johnson replacement collar, Lemma~\ref{lem:C1}, once established,
supplies the actual coefficient equivalence \eqref{eq:HattoriActual}; at
present it is \Open.  The selected class and
divided detector of Section~7 are the geometric instances of the following
pattern.

\subsection{The endpoint module and the Burau action}

\begin{definition}[Endpoint module]\label{def:endpoint-module}
Let $R_\eps=\Z[\eps]/(\eps^7)$, $t=1+\eps$, and
$E_\eps=R_\eps^{88}=\bigoplus_{i=0}^{87}R_\eps e_i$.  One may identify this
free module with the weight-$86$ subspace of the $88$-fold tensor power of
the two-dimensional fundamental module.
\end{definition}

\begin{definition}[Unreduced Burau representation]\label{def:burau}
For the Artin generator $\sigma_i$ of $B_{88}$, with $1\le i\le87$, define
$\rho_t(\sigma_i)$ to be the identity off the span of
$e_{i-1},e_i$ and to have the block
\[
\rho_t(\sigma_i)\big|_{\langle e_{i-1},e_i\rangle}
=
\begin{pmatrix}1-t&t\\1&0\end{pmatrix}
=
\begin{pmatrix}-\eps&1+\eps\\1&0\end{pmatrix}.
\]
Its inverse has block
\[
\rho_t(\sigma_i^{-1})\big|_{\langle e_{i-1},e_i\rangle}
=
\begin{pmatrix}0&1\\t^{-1}&1-t^{-1}\end{pmatrix},
\]
where $t^{-1}=1-\eps+\eps^2-\eps^3+\eps^4-\eps^5+\eps^6$ in $R_\eps$.
For a word $w=\sigma_{i_1}^{s_1}\cdots\sigma_{i_m}^{s_m}$ we use the left
action $\rho_t(w)v=\rho_t(\sigma_{i_1}^{s_1})\cdots
\rho_t(\sigma_{i_m}^{s_m})v$.
\end{definition}

\subsection{From primitive crossings to the cabled word}

The public input stores $252$ primitive crossing rows.  For $i<j$ define the
signed pure-row words
\begin{align}
L_{ij}^{s}
 &=\sigma_{j-1}\cdots\sigma_{i+1}
   \sigma_i^{s}\sigma_i^{s}
   \sigma_{i+1}^{-1}\cdots\sigma_{j-1}^{-1},
\label{eq:L-row}\\
R_{ij}^{s}
 &=\sigma_i\cdots\sigma_{j-2}
   \sigma_{j-1}^{s}\sigma_{j-1}^{s}
   \sigma_{j-2}^{-1}\cdots\sigma_i^{-1}.
\label{eq:R-row}
\end{align}
The $44$-strand word $B_{44}$ is obtained by taking the $R$ rows on segment
$6$ in increasing horizontal coordinate, followed by the $L$ rows on segment
$2$ in decreasing horizontal coordinate.  The cabling homomorphism is
\begin{equation}\label{eq:cabling}
c(\sigma_i)=\sigma_{2i}\sigma_{2i+1}\sigma_{2i-1}\sigma_{2i},
\qquad
c(\sigma_i^{-1})=c(\sigma_i)^{-1},
\end{equation}
and the $88$-strand word is $B_{88}=c(B_{44})$.

\begin{proposition}[Word identities]\label{prop:word-identities}
The exact reconstruction gives
$|B_{44}|=11340$ and $|B_{88}|=45360$.  The Johnson reconstruction verifier
recovers the same $11340$-letter word letter-for-letter.
\end{proposition}

\subsection{The detector and the exact cubic}

Let $R_h=\Q[h]/(h^7)$ and substitute $q=1+h$, $t=q^{-2}$, $\eps=t-1$.
Under the endpoint normalization of the public input, the composite row is
\begin{equation}\label{eq:ell}
\ell=e_{87}^{*}-e_2^{*},
\end{equation}
and the selected shadow vector is
\begin{equation}\label{eq:u}
u=e_2-e_{87}.
\end{equation}
Both formulas use the collar-coordinate endpoint table fixed in the public
input.

\begin{definition}[The divided cubic detector]\label{def:detector}
For a quantum trace class $x$, define
\[
\mathcal D_h(x)=
\widehat C(h)\bigl(\rho_h(W)-I\bigr)P_{86}\Sh_h(x),
\]
where $\Sh_h$ is the quantum trace shadow, $P_{86}$ projects onto the
through-$86$ top cell, and $\widehat C(h)$ is the normalized endpoint cap
row.  When the image lies in $h^3R_h$, the divided cubic detector is
\begin{equation}\label{eq:D3}
D_3(\bar x)=\left.\frac{\mathcal D_h(x)}{h^3}\right|_{h=0}.
\end{equation}
\end{definition}

\begin{proposition}[Selected value]\label{prop:D3-vT}
In $E_\eps$, the first nonzero coefficient of $(\rho_t(B_{88})-I)u$ occurs in
degree $3$.  Contracting with \eqref{eq:ell} gives
\begin{equation}\label{eq:ell-eps-series}
\ell(\rho_t(B_{88})-I)u
=-328\eps^3+14596\eps^4-410246\eps^5+9595271\eps^6.
\end{equation}
Hence, with $[h]\eps=-2$,
\begin{equation}\label{eq:D3-vT-value}
[h^3]\,\ell(\rho_h(B_{88})-I)u=2624\ne0,
\end{equation}
which is $D_3([v_T])$ once Hypothesis~\ref{hyp:P1} identifies $[v_T]$ with
an actual coefficient-trace class.
The complete degree-three vector is printed in Appendix~\ref{app:vector}.
\end{proposition}

\begin{proof}
Traverse $B_{88}$ from right to left, reduce modulo $\eps^7$, and contract
with \eqref{eq:ell}.  The resulting series is
\eqref{eq:ell-eps-series}, so the cubic coefficient is $(-2)^3(-328)=2624$.
\end{proof}

\begin{lemma}[Parameter expansion]\label{lem:eps-of-h}
In $\Z[h]/(h^7)$,
$\eps=(1+h)^{-2}-1=-2h+3h^2-4h^3+5h^4-6h^5+7h^6$.
\end{lemma}

\begin{lemma}[Commutator order]\label{lem:commutator-order}
If $P=I+O(h^a)$ and $Q=I+O(h^b)$ are invertible matrices over an
$h$-adically filtered ring, then $PQP^{-1}Q^{-1}=I+O(h^{a+b})$.  If every
generator of a pure-braid subgroup acts as $I+O(h)$, then its $r$th
lower-central subgroup acts as $I+O(h^r)$.
\end{lemma}

\begin{remark}[Relative class]\label{rem:vT-xi}
The relative class $\xi=\eta_R[T_0]-s_{\mathrm{inv}}\eta_R[T_1]$ satisfies
$D_3(\xi)=0$ because the detector already contains one factor
$\rho_h(W)-I$.  The plain shadow coefficient $[h^3]\ell\Sh_h(\xi)$ equals
$2624$, but that plain shadow is not $D_3(\xi)$.
\end{remark}

\subsection{The grading calculation}

\begin{proposition}[Quantum degree]\label{prop:degree-494}
For the Johnson replacement collar, the selected raw class has quantum
degree $494$.
\end{proposition}

\begin{proof}
There are four contributions:
\begin{align}
d_1&=-44 &&\text{(remove the $44$ Hom normalizations for $N=2$)},\\
d_2&=+227 &&\text{($227$ distinct Frobenius labels)},\\
d_3&=+315 &&\text{($44$ $y$-pairs and $271=44+227$ $z$-pairs)},\\
d_4&=-4 &&\text{(cabled shift for $N=2$, $\alpha=0$, $|r|=2$)}.
\end{align}
Adding gives $-44+227+315-4=494$.  Appendix~\ref{app:grading} collects the
same table without abbreviation.
\end{proof}

\subsection{Quotient algebra and the abstract descent chain}

Let $W_0$ be the one-handle coefficient trace, $W_1$ its quotient by all
two-handle relations, and $W_2$ the quotient of $W_1$ by the three-handle
relations.  Denote the selected class in $W_0$ by $x_0$.

\begin{lemma}[A linear detector survives a quotient]\label{lem:quotient-detector}
Let $V$ be a vector space, $R\subseteq V$ a subspace, and
$\lambda:V\to\Q$ a linear map with $R\subseteq\ker\lambda$.  Then there is a
unique $\bar\lambda:V/R\to\Q$ such that
$\bar\lambda\circ\pi=\lambda$.  If $\lambda(x)\ne0$, then $\pi(x)\ne0$.
\end{lemma}

\begin{proposition}[Abstract nonzero class after quotients]\label{prop:nonzero-W2}
Suppose Hypotheses~\ref{hyp:P1} and~\ref{hyp:P2} hold.  Then there is a
linear functional $\ell_2:W_2\to\Q$ such that
$\ell_2(q_{12}q_{01}x_0)=2624$.  In particular, $q_{12}q_{01}x_0\ne0$.
\end{proposition}

\begin{proof}
Lemma~\ref{lem:HH0-descent} descends the quantum shadow detector through the
coefficient trace.  Hypothesis~\ref{hyp:P1} (status:
Proposition~\ref{thm:Cdischarge}) and Lemma~\ref{lem:quotient-detector}
descend it through $q_{01}$.  Hypothesis~\ref{hyp:P2} (status:
Proposition~\ref{thm:Sdischarge}) and the same quotient lemma descend it
through $q_{12}$.
At every stage the pullback agrees with the preceding row, so
Proposition~\ref{prop:D3-vT} gives the value $2624$.
\end{proof}

\begin{proposition}[Closed candidate class]\label{prop:closed-class}
Suppose Hypotheses~\ref{hyp:P2} and~\ref{hyp:P3} hold and an inhabitant of
Lean \texttt{ExternalGeometry} transports the class through the four-handle
isomorphism.  Then the class of Proposition~\ref{prop:nonzero-W2} maps to a
nonzero element of $\Stwo(X_J)_{0,494}$.  No such inhabitant is constructed.
\end{proposition}

\begin{proof}
MWW Theorems~3.7 and~3.10 identify the three-handle quotient, and
Proposition~3.4 identifies the four-handle stage once cited.
Hypothesis~\ref{hyp:P3} separates the computed empty-link ranks from those
citations.  The Lean implication in Theorem~\ref{thm:A} is conditional on
\texttt{ExternalGeometry}.
\end{proof}

\begin{remark}[Presentation and serialization dependence]\label{rem:serialization}
The raw coefficient $[h^3]\ell(\rho_h(W)-I)u$ depends on the based,
labelled, oriented movie presentation.  Reordering the same crossing rows
by the emitter's increasing-coordinate order produces a different series.
With the corrected collar endpoints its cubic coefficient is $0$ and its first
nonzero coefficient is $663872h^4$.  That order reverses $168$ non-disjoint
events and is not a legal serialization of the Johnson collar.  The former
value $-58976$ used the same mixed endpoint coordinates as the superseded
public value.  Presentation independence must be proved through simultaneous
transport of every part of the naturality square; see
Appendix~\ref{app:legality}.
\end{remark}

\section{A fully worked finite-dimensional example}\label{sec:toy-example}

This section illustrates the algebraic mechanism on a small example.  It is
not a four-manifold calculation and supplies no evidence for any
candidate-specific assumption.

Let $R_0=\Q[\eps]/(\eps^4)$, $t=1+\eps$, and $E_0=R_0^3$ with basis
$e_0,e_1,e_2$.  The two positive Burau generators are
\begin{align}
B_1&=
\begin{pmatrix}
-\eps&1+\eps&0\\
1&0&0\\
0&0&1
\end{pmatrix},
&
B_2&=
\begin{pmatrix}
1&0&0\\
0&-\eps&1+\eps\\
0&1&0
\end{pmatrix}.
\end{align}
Take the pure commutator word
$w_0=\sigma_1^2\sigma_2^2\sigma_1^{-2}\sigma_2^{-2}$.
Multiplying the eight matrices from left to right and reducing modulo
$\eps^4$ gives
\[
\rho(w_0)=
\begin{pmatrix}
1-\eps^3&-\eps^2+\eps^3&\eps^2\\
\eps^2&1&-\eps^2\\
-\eps^2+2\eps^3&\eps^2-3\eps^3&1+\eps^3
\end{pmatrix}.
\]
Let $u_0=e_0-e_1$ and $\ell_0=e_0^*$.  Define
$d_2(v)=[\eps^2]\,\ell_0(\rho(w_0)-I)v$.
Then $d_2=\begin{pmatrix}0&-1&1\end{pmatrix}$ and $d_2(u_0)=1$.
Model two relation maps by the columns
$r_0=(1,0,0)^T$ and $r_1=(0,1,1)^T$.  Then $d_2R=0$, so $d_2$ descends to
$E_0/\im R$ and the selected class is nonzero in the quotient.  What this
example lacks is precisely the geometric content: there is no claim that the
toy relation columns arise from handle maps or that the final quotient is a
smooth invariant.
